\chapter{Conclusions}\label{ch:conclusions}
This thesis presented the design, implementation, and evaluation of a
privacy-preserving decentralised crowdsensing protocol based on smart contracts.
This chapter summarizes the contributions of the work, discusses its
limitations, and outlines directions for future research.

\section{Summary of contributions}\label{sec:conclusion-contributions-summary}
The work presented in this thesis covers three main areas: protocol design,
implementation, and evaluation.

In the design phase, \texttt{PrivacyFirstContract} was developed, building
on the crowdsensing protocol proposed by Bedogni and Ferretti~\cite{bedogni2024crowdsensing}.
A subsequent security and incentive analysis identified several open issues,
which were addressed in \texttt{\seqsplit{PrivacySecondContract}}, an extension that
introduces a geographical reward decay, and a verifier penalty system.

The implementation phase produced two Solidity smart contracts reflecting
the above design, together with a TypeScript SDK layer that exposes a
per-actor API and handles off-chain cryptographic operations, providing a
complete reference implementation of the protocol.

The evaluation phase assessed the contracts along three axes: gas
consumption and its real-world cost across multiple networks, block
capacity contention under concurrent campaign execution, and spatial
coverage and reward distribution on a real-world taxi trajectory dataset
from the city of Porto.

\section{Limitations}\label{sec:conclusion-limitations}
The work presented in this thesis has several limitations.

The geographic simulation was conducted with at most 1\,000 senders due to
RAM constraints on the test machine; a larger population could not be
evaluated without dedicated infrastructure.
Furthermore, the contracts were tested exclusively on a local Hardhat node
and were not deployed on Ethereum mainnet or any L2 network, so the
gas cost figures reported in Section~\ref{sec:gas-analysis} are based on
local measurements rather than live network conditions.

At the protocol level, the centrality of the campaign initiator \texttt{B} remains
an open issue: \texttt{PrivacySecondContract} does not eliminate \texttt{B} as a
trust point, as discussed in Section~\ref{sec:security-analysis-issues}.

Finally, the ProVerif analysis described in Section~\ref{subsec:security-analysis-issues-proverif} assumes
honest enrollment: the sender and verifier enrollment channels are modelled
as private, so the attacker cannot enrol as a malicious participant.
As a consequence, the formal verification only guarantees secrecy properties
under this assumption; scenarios involving a compromised sender or verifier
are outside the scope of the model, as noted in Appendix~\ref{ch:appendix-b-proverif-v1-honest}.

\section{Future work}\label{sec:conclusion-future-work}

Several directions are open for future investigation.
The centrality of the campaign initiator \texttt{B} could be eliminated by replacing
the current key distribution mechanism with Proxy Re-Encryption(PRE)~\cite{ateniese2006pre}: senders
would encrypt their data under their own public key, and a decentralised proxy
network, such as the Threshold Network\cite{thresholdnetwork}, would re-encrypt each ciphertext
for the elected verifier without ever learning the plaintext, removing \texttt{B} as
an online trust point entirely.

The formal security analysis could be strengthened by extending the ProVerif
model to cover dishonest enrollment scenarios, or by porting the verification
to CryptoVerif~\cite{blanchet2008cryptoverif}, which operates in the computational model and provides
guarantees over concrete cryptographic primitives rather than perfect symbolic
abstractions.

On the evaluation side, the geographic simulation should be replicated on
datasets with a more uniform spatial distribution, such as rural or
multi-city traces, to assess the effectiveness of the reward decay mechanism
beyond the concentrated urban setting of the Porto taxi fleet.
Finally, the contracts should be deployed on a public testnet and subsequently
on an L2 network to validate the gas cost figures reported in
Section~\ref{sec:gas-analysis} under real network conditions and to measure
end-to-end latency across all protocol phases.